\chapter{Modèles d'écoulements et similarité dynamique \label{ch:modeles_ecoulement}}

\section{Analyse dimensionnelle}

\subsection{Exemple du pendule simple}
Écartons-nous un instant de la physique des fluides pour nous concentrer sur un exemple simple : un pendule de masse $m$, de longueur $l$, et lâché avec un angle $\alpha$. Les frottements sont négligeables, et le pendule est soumis à la force de gravité. On cherche à déterminer la période $p$ de ses oscillations.

L’un des principes fondamentaux de la physique est que le comportement d’un système ne doit pas dépendre du système d’unités utilisé pour le décrire, car celui-ci est conventionnel. Ainsi, que l’on utilise des grammes, des kilogrammes ou des onces pour mesurer la masse, la période $p$ doit être indépendante de ces choix.

On vérifie cette contrainte en s’assurant de la cohérence dimensionnelle dans la construction des grandeurs physiques. Dans ce problème, nous avons trois unités indépendantes :
\begin{itemize}
  \item une longueur $[l] = L$,
  \item une masse $[m] = M$,
  \item une accélération de gravité $[g] = L/T^2$.
\end{itemize}
On cherche à exprimer un temps $[p] = T$.

Comme la masse $m$ n’intervient pas dans la combinaison dimensionnelle pour $T$, la période ne dépend pas de $m$. La seule combinaison possible pour créer un temps $T$ à partir de $L$ et $L/T^2$ est donc :
\begin{displaymath}
  T = \sqrt{\frac{L}{L/T^2}} = \sqrt{\frac{[l]}{[g]}}.
\end{displaymath}

On en déduit que la période $p$ doit être proportionnelle à $\sqrt{l/g}$ :
\begin{equation}
  p = \sqrt{\frac{l}{g}} \, f(\alpha).
  \label{eq:pendule}
\end{equation}

Nous obtenons ce résultat sans résoudre l’équation de Newton, simplement en imposant la cohérence dimensionnelle. Bien que la fonction $f$ ne soit pas connue, nous pouvons déjà prévoir que pour construire, sur la Lune ($g_l = g/6$), un pendule de même période qu’un pendule de longueur $l$ sur Terre, il faut réduire sa longueur d’un facteur 6. Nous avons aussi prédit que la période d’oscillation du pendule ne dépend pas de sa masse.

\subsection{Théorème de Buckingham et nombres adimensionnels}
Une approche équivalente, bien que plus formelle, consiste à dire que le comportement physique ne peut dépendre que de combinaisons adimensionnelles. Construisons ces nombres pour le pendule. Nous avons cinq grandeurs physiques :
\begin{itemize}
  \item une masse (unité $M$),
  \item une longueur (unité $L$),
  \item une accélération de gravité (unité $L/T^2$),
  \item une période (unité $T$),
  \item un angle $\alpha$ (sans dimension).
\end{itemize}

Nous cherchons à construire une variable adimensionnelle $\Pi$ de la forme :
\begin{displaymath}
  \Pi = m^a l^b g^c p^d \alpha^e.
\end{displaymath}

Pour que $\Pi$ soit adimensionnel, ses exposants doivent annuler les dimensions de $M$, $L$ et $T$. On exprime cette contrainte par les équations suivantes :
\begin{equation}
  \begin{split}
    0 &= a \quad \text{(pour $M$)}, \\
    0 &= b + c \quad \text{(pour $L$)}, \\
    0 &= -2c + d \quad \text{(pour $T$)}.
  \end{split}
  \label{eq:dimensional_system}
\end{equation}

Le système admet pour solutions $a = 0$, $b = -c$, et $d = 2c$. L’exposant $e$ reste indéterminé, car $\alpha$ est déjà adimensionnel. Ce problème est donc décrit par deux nombres adimensionnels indépendants :
\begin{equation}
  \begin{split}
    \Pi^{(1)} &= l^{-1/2} g^{1/2} p, \\
    \Pi^{(2)} &= \alpha.
  \end{split}
  \label{eq:pendule_adim}
\end{equation}

Le nombre de nombres adimensionnels est donné par le \cindex{théorème de Buckingham} : avec 5 quantités physiques ($m$, $l$, $g$, $p$, $\alpha$) et 3 unités indépendantes ($M$, $L$, $T$), on forme $5-3 = 2$ nombres adimensionnels. On peut donc s’attendre à une relation de la forme $F(\Pi^{(1)}, \Pi^{(2)}) = 0$, ou de façon équivalente $\Pi^{(1)} = f(\Pi^{(2)})$. En remplaçant $\Pi^{(1)}$ et $\Pi^{(2)}$ par leurs expressions, on retrouve la relation~\eqref{eq:pendule}.

\subsection{Similarité dynamique et applications}
Voici comment utiliser ce raisonnement en pratique. Supposons que les équations du pendule soient complexes et non résolubles analytiquement, mais que nous souhaitons prédire son comportement. Nous savons que le système est caractérisé par deux nombres adimensionnels. Nous pouvons les interpréter physiquement :
\begin{itemize}
  \item $\Pi^{(1)}$ est lié à la structure du système (longueur, gravité, période).
  \item $\Pi^{(2)}$ dépend de l’angle initial $\alpha$.
\end{itemize}

Si nous voulons prédire la période d’un grand pendule, nous pouvons construire un modèle réduit de même forme, calibrer la fonction $f$ pour différentes valeurs de $\alpha$, puis appliquer la relation~\eqref{eq:pendule} au grand pendule. On peut par exemple s'attendre à ce que $\Pi^{(1)}$ dépende de la façon dont la masse est répartie le long de l'axe du pendule, dont nous n'avons toujours rien dit. 

On dit que deux systèmes ayant les mêmes nombres adimensionnels sont dans une relation de \cindex{similarité dynamique}. Cette notion est fondamentale pour prédire le comportement de systèmes physiques complexes (comme un avion) à partir de tests en soufflerie avec des modèles réduits.

\subsection{Prise en compte de la rotation terrestre}
Supposons maintenant que la vitesse angulaire de rotation $\Omega$ de la Terre soit pertinente. C’est une nouvelle variable, mais elle n’introduit pas de nouvelle dimension (son unité est $T^{-1}$). Nous créons donc une nouvelle variable adimensionnelle $\Pi^{(3\star)} = p\Omega$. Cependant, comme $p$ apparaît déjà dans $\Pi^{(1)}$, il est judicieux de construire une variable indépendante de $p$ :

\begin{displaymath}
  \Pi^{(3)} = \frac{\Pi^{(3\star)}}{\Pi^{(1)}} = \Omega \sqrt{\frac{l}{g}}.
\end{displaymath}

Notre nouvelle relation devient :
\begin{equation}
  p = \sqrt{\frac{l}{g}} \, f\left(\Omega \sqrt{\frac{l}{g}}, \alpha\right).
  \label{eq:grand_pendule}
\end{equation}

En pratique, on constate presque toujours que dès qu’un nombre adimensionnel est soit très petit (disons, $\ll 0.01$) ou très grand ($\gg 100$), la valeur exacte qu’il prend perd son importance. C’est le fait qu’il soit très petit ou très grand qui est pertinent. Pour un petit pendule, tel qu’on pourrait le construire dans une salle de classe, $\Pi^{(3)} \sim 2 \cdot 10^{-5}$. On peut en déduire que si la vitesse de rotation de la Terre était deux fois plus rapide, l’expérience ne serait pas fortement modifiée. De même, on prédit que pour $\alpha$ très petit, la valeur précise de $\alpha$ perd son importance. 

On peut aussi estimer la longueur typique d’un pendule compatible avec des effets visibles de la rotation terrestre, c’est-à-dire la longueur qu’il faudrait pour mesurer commodément $\Omega$. En imposant $\Pi^{(3)} = 1$, on trouve $l \approx 1800~\textrm{km}$. C’est la longueur typique des \cindex{oscillations d’inertie} visibles dans l’océan.

Dans ce raisonnement, nous avons utilisé des dimensions typiques ($l \sim 1~\textrm{m}$, $g = 9.81~\textrm{m}/\textrm{s}^2$, etc.) pour prédire l’ordre de grandeur de $p$ ou l’importance relative d’effets physiques. C’est le cœur de l’analyse dimensionnelle.

Notons enfin qu’il est souvent plus pertinent d’utiliser des vitesses angulaires que des périodes ou des fréquences. Par exemple, pour un oscillateur harmonique de pulsation $\omega$, l’équation du mouvement s’écrit $\ddot{x} + \omega^2 x = 0$. La dérivée temporelle fait apparaître $\omega$ naturellement, ce qui rend son utilisation plus directe que celle de la période $T = 2\pi/\omega$. 

D'ailleurs, pour un pendule de masse ponctuelle, et en utiliant la formule \eqref{eq:pendule_adim}, la théorie Newtonienne pour des angles $\alpha$ petits prédit $\Pi^{(1)}=2\pi$,  alors qu'on définissant, de façon plus astucieuse, $\Pi^{(1)}=l^{-1/2}g^{1/2}\omega^{-1}$, on a $\Pi^{(1)}$.


\section{Analyse dimensionnelle de l'équation de Navier-Stokes}

\subsection{Transition vers le régime turbulent \label{ssec:turbulent}}

Reprenons l'équation décrivant l'écoulement dans le plan de Couette \eqref{eq:couette_plan}. Nous avons expliqué  que le terme de transport de moment cinétique s'annule \emph{si la symmétrie des conditions frontières se reflète dans le champ de vitesse}. Rien ne le garantit a priori cependant. Il est possible en principe que l'écoulement que nous avons déterminé se révèle instable. Les conditions de stabilité sont cependant très difficiles à établir analytiquement. C'est un problème de stabilité hydrodynamique qui peut en principe se réduire à sa dimension mathématique: quelles sont les conditions d'instabilité ? Ces conditions ne peuvent pas dépendre du système d'unité utilisé pour décrire l'écoulement. 

Reprenons donc l'équation \eqref{eq:couette_plan} mais en utilisant la \cindex{viscosité cinématique} $\nu = \mu/\rho$, qui a pour avantage de s'exprimer en $\mathrm{m}^2/s$ et nous débarasser des kilogrammes:

\begin{equation}
u \od{u}{x} = -{C} + \nu \od[2]{u}{x}. 
\label{eq:couette_cinematique}
\end{equation}

Le problème, avec ses conditions frontières, met en scène 3 grandeurs : $u$, $h$ (la distance entre les deux plaques), $\nu$, et deux unité indépendantes: l'unité de longueur, exprimée en $m$, et de temps $s$ (les unités de vitesse et la viscosité cinématique sont construites à partir de ces deux unités). La structure de la solution va donc être décrite par \emph{un} nombre adimensionel. 

Il y a \textit{a priori} différentes façons de le construire, mais nous pouvons en chercher une qui a un sens physique bien défini. Reprenons l'équation \eqref{eq:couette_cinematique}. Le membre de gauche (l'advection) a les unités $U\cdot U/L$, où $U$ symbolise l'unité de vitesse, et $L$, d'espace. Cela signifie que si l'on change les dimensions du système, ce terme va rester proportionnel à $u^2/h$. 
Le membre de droite, qui exprime les forces de viscosité,  est lui proportionnel à $\nu u^2/h$. Nous pouvons former un nombre adimensionnel en mettant en rapport le terme d'advection avec celui faisant intervenir la viscosité. Appelons-le le nombre de Reynolds: 

\begin{equation}
  \Rey \eqdef \frac{u^2}{h} \cdot \frac{h^2}{\nu u} = \frac{uh}{\nu}. 
  \label{eq:reynolds_nr}
\end{equation}



Les principes de l'analyse dimensionnelle nous garantissent que la condition de stabilité de l'écoulement turbulent ne dépend \emph{que} de ce nombre, pour autant que nous ayons correctement identifié toutes les contraintes physiques du problème. Par ailleurs, notre intuition physique nous permet de prédire que l'écoulement laminaire perdra sa stabilité lorsque les termes d'advection deviennent potentiellement plus importants, suggérant une transition vers en \cindex{écoulement turbulent} \emph{au delà} d'un nombre de Reynolds critique. Ce nombre de Reynolds peut être déterminé par une expérience réalisée par Osborne Reyolds, illustrée p. 71 (Fig. 15) du Vol. II des "Papers on mechanical and physical subjects" \cite{reynolds01aa}.

Dans une conduite cylindrique, la transition s'opère pour des nombre de Reynolds entre 2000 et 3000.  Cette mesure donne lieu à une classification des écoulements qui s'étend bien au delà des conditions particulières de l'écoulement de Poiseuille. 

Ainsi, pour un écoulement caractérisé par une vitesse typique $U$, une dimension typique $L$ et une viscosité cinématique $\nu$, on définit $\Rey = UL/\nu$ et on considère:

\begin{itemize}
  \item $\Rey \lesssim 1$ : les termes d'advection sont dominés par les termes de viscosité. Ce sont les conditions de l'écoulement de Stokes tel que nous l'avons dérivé en section \ref{eq:stokes} et dont nous avons fait l'hypothèse pour le \cindex{problème de sédimentation}
  \item $5 \lesssim \Rey \lesssim 1000 $: l'écoulement est principalement laminaire mais les forces et dissipations mesurées commencent à s'écarter des conditions de l'écoulement de Stokes. Des tourbillons apparaîssent aux points d'asymmétrie imposés par l'experience, par exemple à l'embouchure du tube, ou, si on considère l'écoulement autour d'une obstacle, à l'arrière de l'obstacle. 
  \item $2000 \lesssim \Rey \lesssim 3000$ : régime intermittent ou des phases turbulentes succèdent à des phases laminaires;
  \item $3000 \lesssim \Rey$ : régime turbulent, dont la nature et la complexité peuvent évoluer en fonction de $\Rey$. 
\end{itemize}

Rappelons que le seuil critique $\Rey \sim 3000$ s'applique à une expérience bien précise. Le problème d'ingéniérie posé aux concepteurs d'engins aéorodynamiques est précisément de concevoir des profils et des matérieux qui font reculer la transition turbulence vers des nombres de Reynolds plus élevés.  

\subsection{Prédiction de l'épaisseur de la couche limite}

Considérons un écoulement autour d’un obstacle de longueur caractéristique $L$. Près de la paroi, la vitesse du fluide doit s’annuler (condition de non-glissement). La viscosité y joue donc un rôle dominant, même si $\Rey \gg 1$ globalement.

Nous cherchons l’épaisseur $\delta$ de la \cindex{couche limite}, où les effets visqueux sont significatifs. Dans cette région, les termes inertiels et visqueux dans l’équation de Navier-Stokes doivent être du même ordre de grandeur. Ainsi, par rapport au cadre posé dans la section \ref{ssec:turbulent}, où nous cherchions un nombre de Reynolds critique pour la transition vers un régime turbulent, ici nous cherchons la dimension typique pour laquelle les termes inertiels et visqueux peuvent s'équilibrer, étant donné un nombre de Reynolds imposé par l'expérience. 

Analysons les les ordres de grandeurs présents dans l'équation de Navier-Stokes. Si $U$ est une vitesse typique, et $L$ une distance caractéristique de l'objet aérodynamique étudié, le terme d'advection sera de l'ordre de $U^2/L$. Dans la couche limite, le terme de viscosité met en scène une dissipation visqueuse sur la distance typique de la couche limite, donc un terme de l'ordre de $\nu U/\delta^2$. La condition est donc: 

\begin{displaymath}
  \frac{U^2}{L} \sim \nu \frac{U}{\delta^2},\textrm{ soit} \delta^2 \sim \frac{L^2}{\Rey}. 
\end{displaymath}

Ainsi, l’analyse dimensionnelle prédit que, pour $L$ donné (par la dimension de l'obstacle), l’épaisseur de la couche limite $\delta$ décroît comme $1/\sqrt{\Rey}$. Cette prédiction explique pourquoi, à haut nombre de Reynolds, la viscosité ne joue un rôle significatif que dans une fine couche près des parois, tandis que l’écoulement est essentiellement non visqueux ailleurs. 


\begin{redmargin}
  Nous avons pris un risque. Nous avons dérivé le terme visqueux en supposant qu'il a la forme donnée dans l'équation \eqref{eq:ns_incompressible} et en estimant que le Laplacien aura bien un ordre de grandeur proportionnel à $\delta^{-2}$. A vrai dire nous n'en savons rien car l'écoulement peut prendre une forme complexe et multidimensionnelle. 

  La seule chose que nous pouvons dire est ceci. Notre écoulement est maintenant caractérisé par 4 gradeurs imposées par l'expérience : $U$, $L$, $\delta$ et $\nu$, et deux dimesions. Il y a donc deux nombres adimensionnels, que l'on peut construire comme $\Rey$ et $\delta/L$. La physique impose une relation de type $F(\delta/L, \Rey)=0$, dont la résolution peut prendre la forme $\delta/L = f(\Rey)$, soit $\delta = L f(\Rey)$. 

  Nous ne connaissons pas à priori la relation $f(\Rey)$ mais nous estimons qu'à la limite d'un régime laminaire il doit prendre la forme $f(\Rey)\propto \Rey^{-0.5}$. Dans un régime plus complexe, on attendra plus généralement $f(\Rey) \propto \Rey^-\alpha$ et, généralement, $\alpha<0.5$, avec une valeur typique de l'ordre de 0.2. 
\end{redmargin}


\subsection{Discussion} 

\begin{studentinvestigation}{Simulateur de vol}
  Soit un avion de $100,\mathrm{m}$, une vitesse de $300\,\mathrm{m}/{s}$, une viscosité cinématique de l'air de $1.5\ 10^{-5}\mathrm{m^2}/\mathrm{s}$. Quel vaut $\Rey$ ? Quel est l'ordre de grandeur de  $\delta$ dans l'hypothèse laminaire ? 

  Sur cette base, supposons que l'on désire calculer l'écoulement  avec un superordinateur avec la résolution nécessaire pour capture la couche limite. De combien de maille tri-dimensionnelles avons-nous besoin ? 
\end{studentinvestigation}

Si nous tenons compte du critère dit CFL (vu au cours LPHYS1303), nous pouvons aussi calculer la résolution temporelle à laquele l'ordinateur doit calculer. On trouvera .... 20 micro-secondes ! On s'apercevra très vite que résoudre les équations de Navier-Stokes autour d'un avion est impossible. Il faudrait de l'ordre de $10^{16}$ opérations pour simuler une seconde de vol.  Le même constat s'applique aux calcul de flux atmosphériques et océaniques. Les équations de Navier-Stokes doivent  donc être adaptées (modifiées ou tronquées) pour éviter de résoudre dans le détail les écoulements aux échelles les plus petites, et qui, pourtant interviennent dans les bilans d'énergie et de quantité de mouvement. Cela reste une des grands enjeux de la mécanique des fluides contemporaine. 

%
% L’analyse dimensionnelle permet de :
% \begin{itemize}
%   \item Comprendre pourquoi les écoulements à grand $\Rey$ (avions, navires) sont souvent modulés par des couches limites minces.
%   \item Expliquer l’importance de la rugosité des surfaces (même une petite rugosité peut perturber la couche limite si $\delta$ est très fin).
%   \item Prédire l’apparition de la turbulence lorsque les perturbations dans la couche limite deviennent instables (pour $\Rey$ critique dépendant de la géométrie).
% \end{itemize}
%
% Cependant, cette analyse reste qualitative. Pour une description quantitative, il faut résoudre les équations de la couche limite (équations de Prandtl), qui sortent du cadre de ce cours.
%
% \subsection{Limites des simulations numériques directes (DNS)}
% L’analyse dimensionnelle permet non seulement de comprendre les régimes d’écoulement, mais aussi d’estimer les ressources computationnelles nécessaires pour les simuler numériquement. Considérons un écoulement autour d’un obstacle (comme un avion) à grand nombre de Reynolds $\Rey$.
%
% \subsubsection{Échelle de la grille et nombre de mailles}
% Pour résoudre les équations de Navier-Stokes sans modèle de turbulence (simulation directe, ou DNS), il faut capturer toutes les échelles de l’écoulement, y compris les plus petites, où la dissipation visqueuse domine. Dans la couche limite, l’épaisseur $\delta$ des structures visqueuses est donnée par :
% \begin{displaymath}
%   \delta \sim \frac{L}{\sqrt{\Rey}},
% \end{displaymath}
% où $L$ est la taille caractéristique de l’obstacle (ex. : envergure d’une aile).
%
% Pour résoudre ces structures, la taille des mailles $\Delta x$ de la grille numérique doit être de l’ordre de $\delta$ :
% \begin{displaymath}
%   \Delta x \sim \delta \sim \frac{L}{\sqrt{\Rey}}.
% \end{displaymath}
%
% Dans un domaine 3D de taille $L^3$, le nombre total de mailles $N$ est donc :
% \begin{displaymath}
%   N \sim \left(\frac{L}{\Delta x}\right)^3 \sim \left(\sqrt{\Rey}\right)^3 = \Rey^{3/2}.
% \end{displaymath}
%
% \subsubsection{Coût temporel et critère CFL}
% Le pas de temps $\Delta t$ d’une simulation explicite est contraint par le **critère de stabilité CFL** (Courant-Friedrichs-Lewy), qui impose :
% \begin{displaymath}
%   \Delta t \lesssim \frac{\Delta x}{U},
% \end{displaymath}
% où $U$ est la vitesse caractéristique de l’écoulement. Le nombre total de pas de temps $N_t$ pour simuler un temps physique $T$ est donc :
% \begin{displaymath}
%   N_t \sim \frac{T}{\Delta t} \sim \frac{T U}{\Delta x} \sim T U \sqrt{\Rey}.
% \end{displaymath}
%
% Le coût computationnel total $C$ (proportionnel à $N \times N_t$) croît ainsi comme :
% \begin{displaymath}
%   C \sim \Rey^{3/2} \times \Rey^{1/2} = \Rey^2.
% \end{displaymath}
%
% \subsubsection{Implications pratiques}
% Pour des écoulements réalistes (ex. : avion en vol, $\Rey \sim 10^7$–$10^9$), une DNS devient **prohibitive** :
% \begin{itemize}
%   \item Pour $\Rey = 10^6$, $N \sim 10^{9}$ mailles et $N_t \sim 10^6$ pas de temps, soit un coût $\sim 10^{15}$ opérations.
%   \item Les supercalculateurs modernes (ex. : $\sim 10^{18}$ op/s) mettraient des **jours** à simuler quelques secondes physiques.
% \end{itemize}
%
% \begin{key}{Conséquences pour la modélisation}
% Cette analyse justifie l’utilisation de modèles de turbulence (LES, RANS) pour les écoulements à grand $\Rey$ :
% \begin{itemize}
%   \item Les **Large Eddy Simulations (LES)** résolvent les grandes échelles et modélisent les petites.
%   \item Les **Reynolds-Averaged Navier-Stokes (RANS)** moyennent les équations et modélisent entièrement la turbulence.
% \end{itemize}
% Ces approches réduisent drastiquement le coût computationnel, au prix d’une perte de précision sur les petites échelles.
% \end{key}
%
% \subsubsection{Exemple : Simulation d’un avion}
% Pour un avion de taille $L \sim 10~\textrm{m}$ volant à $U \sim 250~\textrm{m/s}$ dans l’air ($\nu \sim 1.5 \cdot 10^{-5}~\textrm{m}^2/\textrm{s}$), le nombre de Reynolds est :
% \begin{displaymath}
%   \Rey = \frac{UL}{\nu} \sim \frac{250 \times 10}{1.5 \cdot 10^{-5}} \sim 1.7 \cdot 10^8.
% \end{displaymath}
% Une DNS nécessiterait donc :
% \begin{itemize}
%   \item $\sim (1.7 \cdot 10^8)^{3/2} \sim 10^{12}$ mailles,
%   \item $\sim 10^{12} \times 10^4 = 10^{16}$ opérations pour simuler 1 seconde de vol.
% \end{itemize}
% C’est hors de portée des capacités actuelles, même avec les supercalculateurs.

